\documentclass[12pt]{article}

\input{~/School/Notes/header.tex}

\usepackage{booktabs}

\setlength{\parindent}{1.25cm}
\setlength{\parskip}{0.25em}
\renewcommand{\arraystretch}{1.15}

\title{Лабораторная работа по физике}
\author{Денисов Илья}
\date{}

\begin{document}

    \maketitle

    \vspace{5cm}
    {\centering \LARGE Серый ящик с тремя выводами \par}

    \newpage

    \tableofcontents
    \setcounter{tocdepth}{3}

    \newpage

    \section{Цель работы}

    Установить, по какой схеме (<<звезда>> или <<треугольник>>) и из каких элементов (резисторы или конденсаторы) собран серый ящик с тремя выводами, и определить сопротивления его резисторов и ёмкость конденсатора.

    \section{Задачи}

    \begin{enumerate}
        \item Измерить омметром сопротивление между каждой парой выводов $1$--$2$, $1$--$3$, $2$--$3$ и по характеру показаний определить схему и тип элементов.
        \item Найти сумму сопротивлений резисторов $R_1 + R_2$.
        \item Измерить напряжения $U'_2$ и $U'_1$, найти отношение $R_2/R_1$ и сопротивления $R_1$ и $R_2$.
        \item Измерить напряжения $U_3$ и $U_4$ без дополнительного конденсатора и с конденсатором известной ёмкости $C_0$ и найти ёмкость $C$.
        \item Оценить случайные погрешности результатов.
    \end{enumerate}

    \section{Оборудование и метод измерения}

    Использовались серый ящик с тремя пронумерованными выводами, мультиметр (в режимах омметра и вольтметра), конденсатор известной ёмкости $C_0 = 10$ мкФ и соединительные провода.

    Омметр представляет собой источник постоянной ЭДС $\mathcal{E}$ с некоторым внутренним сопротивлением $r$, поэтому его можно использовать как источник напряжения для зарядки конденсатора. Мультиметр в режиме вольтметра имеет большое сопротивление $R_V$, которое в работе не измерялось и в расчёты численно не входит. О том, что оно велико, говорит медленное убывание показаний при разрядке конденсатора на вольтметр. Фиксируя второе-третье показание вольтметра, можно сравнительно точно найти напряжение в начальный момент.

    \subsection{Определение схемы и типа элементов}

    Мультиметр в режиме омметра подключался попарно к выводам $1$--$2$, $1$--$3$ и $2$--$3$. Результаты наблюдений приведены в таблице.

    \begin{table}[H]
        \centering
        \begin{xltabular}{\textwidth}{@{}lXX@{}}
            \toprule
            Выводы & Показания омметра & Характер показаний\\
            \midrule
            $1$--$2$ & $114$ кОм & устойчивое значение\\
            $1$--$3$ & от конечного значения до $\infty$ & растут за несколько секунд\\
            $2$--$3$ & от конечного значения до $\infty$ & растут за несколько секунд\\
            \bottomrule
        \end{xltabular}
        \caption{Показания омметра для разных пар выводов}
        \label{tab:ohm}
    \end{table}

    Рост показаний до бесконечности означает, что идёт зарядка конденсатора: в установившемся режиме ток через него не течёт, а между выводами $1$--$3$ и $2$--$3$ нет чисто резистивного пути. Устойчивое конечное значение между выводами $1$--$2$ возможно, только если на пути между ними нет конденсаторов.

    \textbf{Треугольник невозможен.} При соединении <<треугольником>> между любыми двумя выводами есть два параллельных пути: прямой элемент и цепочка из двух остальных. Бесконечное сопротивление между выводами $1$--$3$ и $2$--$3$ достигается, только если оба элемента, подключённых к выводу $3$, являются конденсаторами. Тогда между выводами $1$--$2$ параллельно третьему элементу стоит цепочка из этих двух конденсаторов, через которую в начале идёт ток зарядки, и показания менялись бы со временем (а при конденсаторе в качестве третьего элемента они были бы бесконечными). Это противоречит устойчивому значению $114$ кОм.

    \textbf{Звезда.} В схеме типа <<звезда>> сопротивление между выводами $1$--$2$ равно сумме двух элементов и остаётся постоянным, поэтому эти элементы~--- резисторы $R_1$ (у вывода $1$) и $R_2$ (у вывода $2$). Показания между выводами $1$--$3$ и $2$--$3$ растут до бесконечности, значит, элемент у вывода $3$~--- конденсатор $C$. Схема ящика показана на рисунке.

    \begin{figure}[H]
        \centering
        \begin{tikzpicture}[scale=1, line join=round, line cap=round]
            % Центральный узел
            \fill (0, 0) circle (2pt);
            % Резистор R_1 (вывод 1)
            \draw (0, 0) -- (0, 0.6);
            \draw (-0.2, 0.6) rectangle (0.2, 1.8);
            \draw (0, 1.8) -- (0, 2.3);
            \draw[fill=white] (0, 2.4) circle (0.08);
            \node[above] at (0, 2.5) {$1$};
            \node[right] at (0.3, 1.2) {$R_1$};
            % Резистор R_2 (вывод 2)
            \begin{scope}[rotate=120]
                \draw (0, 0) -- (0, 0.6);
                \draw (-0.2, 0.6) rectangle (0.2, 1.8);
                \draw (0, 1.8) -- (0, 2.3);
                \draw[fill=white] (0, 2.4) circle (0.08);
            \end{scope}
            \node[below left] at (210:2.5) {$2$};
            \node at ($(210:1.2) + (120:0.55)$) {$R_2$};
            % Конденсатор C (вывод 3)
            \begin{scope}[rotate=-120]
                \draw (0, 0) -- (0, 0.9);
                \draw[thick] (-0.3, 0.9) -- (0.3, 0.9);
                \draw[thick] (-0.3, 1.1) -- (0.3, 1.1);
                \draw (0, 1.1) -- (0, 2.3);
                \draw[fill=white] (0, 2.4) circle (0.08);
            \end{scope}
            \node[below right] at (330:2.5) {$3$};
            \node at ($(330:1.0) + (60:0.55)$) {$C$};
        \end{tikzpicture}
        \caption{Схема серого ящика (<<звезда>>) с номерами выводов}
        \label{fig:star}
    \end{figure}

    \section{Рабочие формулы}

    Сумма сопротивлений резисторов найдена омметром между выводами $1$--$2$:
    \[
        R_1 + R_2 = 114 \ \text{кОм}.
    \]

    \subsection{Сопротивления резисторов}

    \textbf{Измерение $U'_2$.} Омметр подключается к выводу $1$ и к закороченным выводам $2$ и $3$. Конденсатор оказывается включённым параллельно $R_2$ и заряжается до напряжения на этом резисторе:
    \[
        U_2 = \frac{\mathcal{E} R_2}{R_1 + R_2 + r}.
    \]
    Затем омметр отключается, перемычка между выводами $2$ и $3$ снимается, и вольтметр подключается к выводам $1$ и $3$: получается цепь из последовательно соединённых $R_1$, $C$ и вольтметра (вывод $2$ свободен, ток через $R_2$ не течёт). В начальный момент
    \[
        U'_2 = U_2 \cdot \frac{R_V}{R_1 + R_V} = \frac{\mathcal{E} R_2}{R_1 + R_2 + r} \cdot \frac{R_V}{R_1 + R_V}.
    \]

    \textbf{Измерение $U'_1$.} Аналогично омметр подключается к выводу $2$ и к закороченным выводам $1$ и $3$, а после снятия перемычки вольтметр подключается к выводам $2$ и $3$:
    \[
        U_1 = \frac{\mathcal{E} R_1}{R_1 + R_2 + r}, \qquad
        U'_1 = U_1 \cdot \frac{R_V}{R_2 + R_V} = \frac{\mathcal{E} R_1}{R_1 + R_2 + r} \cdot \frac{R_V}{R_2 + R_V}.
    \]

    Делим $U'_2$ на $U'_1$; неизвестные $\mathcal{E}$ и $r$ сокращаются:
    \[
        \frac{U'_2}{U'_1} = \frac{R_2}{R_1} \cdot \frac{R_2 + R_V}{R_1 + R_V} \approx \frac{R_2}{R_1},
    \]
    так как $R_1, R_2 \ll R_V$. Относительная поправка равна $|R_2 - R_1| / R_V \approx 24 \ \text{кОм} / R_V$ и мала при $R_V$ порядка мегаом и выше, поэтому далее она не учитывается. Пусть $k = U'_2 / U'_1$. Тогда
    \[
        R_1 = \frac{R_1 + R_2}{1 + k}, \qquad R_2 = (R_1 + R_2) - R_1.
    \]

    \subsection{Ёмкость конденсатора}

    \textbf{Измерение $U_3$.} Омметр подключается к выводам $2$ и $3$ на несколько минут; конденсатор заряжается до $\mathcal{E}$, так как в установившемся режиме ток через $R_2$ и $r$ не течёт. Затем вольтметр подключается к выводам $2$ и $3$, и в начальный момент
    \[
        U_3 = \frac{R_V}{R_2 + R_V}\,\mathcal{E}.
    \]
    Перед каждым новым измерением конденсатор разряжается замыканием выводов $3$ и $1$.

    \textbf{Измерение $U_4$.} К выводу $3$ подключается конденсатор известной ёмкости $C_0$, и цепочка $C$, $C_0$ заряжается от омметра через $R_2$ в течение нескольких минут. Установившееся напряжение на конденсаторе $C$ равно $\mathcal{E} C_0 / (C_0 + C)$. После отключения $C_0$ напряжение между выводами $2$ и $3$ в начальный момент равно
    \[
        U_4 = \frac{R_V}{R_2 + R_V} \cdot \frac{C_0}{C_0 + C}\,\mathcal{E}.
    \]
    Перед каждым измерением оба конденсатора разряжаются.

    Делим $U_3$ на $U_4$; множители с $R_V$, $R_2$ и $\mathcal{E}$ сокращаются:
    \[
        \frac{U_3}{U_4} = \frac{C_0 + C}{C_0}, \qquad C = C_0\left(\frac{U_3}{U_4} - 1\right).
    \]

    \section{Результаты измерений}

    Напряжения измерены в вольтах; в таблице приведены результаты каждой серии и их средние значения.

    \begin{table}[H]
        \centering
        \begin{xltabular}{\textwidth}{@{}lXXXXXX@{}}
            \toprule
            $N$ & 1 & 2 & 3 & 4 & 5 & Среднее\\
            \midrule
            $U'_2$, В & 0.0532 & 0.0540 & 0.0536 & 0.0543 & 0.0538 & 0.05378\\
            $U'_1$, В & 0.0811 & 0.0828 & 0.0822 & 0.0820 & 0.0821 & 0.08204\\
            $U_3$, В  & 0.264  & 0.265  & 0.267  & 0.265  & 0.266  & 0.2654\\
            $U_4$, В  & 0.247  & 0.250  & 0.248  & 0.249  & 0.248  & 0.2484\\
            \bottomrule
        \end{xltabular}
        \caption{Результаты серий измерений напряжений}
        \label{tab:data}
    \end{table}

    \section{Обработка результатов}

    \subsection{Сопротивления резисторов}

    Отношение напряжений:
    \[
        k = \frac{\langle U'_2 \rangle}{\langle U'_1 \rangle} = \frac{0.05378}{0.08204} = 0.6555 = \frac{R_2}{R_1}.
    \]
    Так как $R_1 + R_2 = 114$ кОм, то
    \[
        R_1 = \frac{114}{1 + 0.6555} = 68.9 \ \text{кОм}, \qquad R_2 = 114 - 68.9 = 45.1 \ \text{кОм}.
    \]

    \subsection{Ёмкость конденсатора}

    \[
        \frac{\langle U_3 \rangle}{\langle U_4 \rangle} = \frac{0.2654}{0.2484} = 1.068, \qquad
        C = 10 \cdot (1.068 - 1) = 0.68 \ \text{мкФ}.
    \]

    \subsection{Оценка случайных погрешностей}

    Для каждой серии из $n = 5$ измерений находим среднеквадратичную погрешность среднего и случайную погрешность при доверительной вероятности $\alpha = 0.90$:
    \[
        \sigma_{\bar U} = \sqrt{\frac{\sum_i (U_i - \bar U)^2}{n(n - 1)}}, \qquad
        \varDelta U = t_{0.90;\,4}\,\sigma_{\bar U}, \qquad t_{0.90;\,4} = 2.132.
    \]

    \begin{table}[H]
        \centering
        \begin{xltabular}{\textwidth}{@{}lXXX@{}}
            \toprule
            Величина & $\sigma_{\bar U}$, мВ & $\varDelta U$, мВ & $\varepsilon_U$, \%\\
            \midrule
            $U'_2$ & 0.19 & 0.40 & 0.74 \\
            $U'_1$ & 0.27 & 0.58 & 0.71 \\
            $U_3$  & 0.51 & 1.09 & 0.41 \\
            $U_4$  & 0.51 & 1.09 & 0.44 \\
            \bottomrule
        \end{xltabular}
        \caption{Случайные погрешности средних значений ($\alpha = 0.90$)}
        \label{tab:err}
    \end{table}

    Для отношения $k$:
    \[
        \varepsilon_k = \sqrt{\varepsilon_{U'_2}^2 + \varepsilon_{U'_1}^2} = 1.0\%, \qquad \varDelta k = 0.6555 \cdot 0.010 = 0.0067,
    \]
    \[
        \varDelta R_1 = \varDelta R_2 = \frac{(R_1 + R_2)\,\varDelta k}{(1 + k)^2} = \frac{114 \cdot 0.0067}{1.6555^2} = 0.28 \ \text{кОм}.
    \]
    Для ёмкости, где $q = U_3 / U_4$:
    \[
        \varepsilon_{q} = \sqrt{\varepsilon_{U_3}^2 + \varepsilon_{U_4}^2} = 0.60\%, \qquad
        \varDelta q = 1.068 \cdot 0.0060 = 0.0064, \qquad \varDelta C = C_0\,\varDelta q = 0.064 \ \text{мкФ}.
    \]
    Оценка учитывает только случайный разброс серий; значения $R_1 + R_2$ и $C_0$ приняты точными, собственная погрешность мультиметра не учитывалась.

    \section{Окончательные результаты}

    С доверительной вероятностью $\alpha = 0.90$:
    \[
        R_1 = (68.9 \pm 0.3) \ \text{кОм}, \qquad R_2 = (45.1 \pm 0.3) \ \text{кОм}, \qquad C = (0.68 \pm 0.06) \ \text{мкФ}.
    \]

    \section{Вывод}

    \begin{enumerate}
        \item Серый ящик собран по схеме <<звезда>>. Между выводами $1$--$2$ омметр показывает устойчивое значение $114$ кОм, значит, элементы у выводов $1$ и $2$~--- резисторы $R_1$ и $R_2$. Между выводами $1$--$3$ и $2$--$3$ показания растут до бесконечности, то есть идёт зарядка конденсатора $C$, который подключён к выводу $3$. Схема <<треугольник>> не подходит, так как при ней показания между выводами $1$--$2$ тоже менялись бы со временем.

        \item Сумма сопротивлений резисторов $R_1 + R_2 = 114$ кОм. Из отношения $U'_2 / U'_1 = R_2 / R_1 = 0.6555$ получены $R_1 = (68.9 \pm 0.3)$ кОм и $R_2 = (45.1 \pm 0.3)$ кОм. Относительная погрешность отношения составляет около $1\%$, поэтому сопротивления определены с точностью до долей процента. Поправка на конечное сопротивление вольтметра не учитывалась: она мала при $R_V \gg R_1, R_2$, а при вычислении $C$ сокращается точно.

        \item Ёмкость конденсатора $C = (0.68 \pm 0.06)$ мкФ. Её относительная погрешность ($\approx 9\%$) заметно больше, чем у остальных величин, так как $C \ll C_0$: отношение $U_3 / U_4 = 1.068$ мало отличается от единицы, и малая погрешность напряжений сильно сказывается на разности $U_3 / U_4 - 1$. Точность была бы выше при $C_0$, сравнимом с $C$.

        \item Для повышения точности использовалось второе-третье показание вольтметра (до заметной разрядки конденсатора через вольтметр), несколько минут ожидания до полной зарядки и разрядка конденсаторов перед каждым измерением. Напряжения в каждой паре ($U'_2$ и $U'_1$, $U_3$ и $U_4$) снимаются в одинаковый момент после подключения вольтметра, а разрядка идёт практически с одной и той же постоянной времени, поэтому она почти одинаково уменьшает оба показания и на их отношение влияет слабо. Каждая серия содержит $5$ измерений; точность можно повысить, увеличив их число.
    \end{enumerate}

\end{document}
